%------------------------------------------------------------------------------%
\section{Função Logarítmica}

Se considerarmos a função exponencial
$f\colon \R \to \R_+^*$,
definida por $f(x) = a^x$, $a>0$ e $a \neq 1$, então temos
que $f$ é uma bijeção sobre a sua imagem, portanto possui inversa.
Sua inversa é denominada
\emph{função logarítmica na base $a$}\index{Função!logaritmica},
$a > 0$ e $a \neq 1$.
\begin{align*}
  \log_a\colon \R_{+}^{*} & \to \R        \\
  x                       & \to \log_a x
\end{align*}
e é definida por
\[
  \log_a x = y
  \quad\Leftrightarrow\quad
  a^y = x.
\]

Vamos fazer alguns exemplos de cálculo de logaritmos,

%------------------------------------------------------------------------------%
\begin{example}
  Encontre $\log_2 8$:

  \solution
  Usando a definição de logaritmo temos
  \begin{align*}
    \log_2 8 & = x \\
    2^x      & = 8 \\
    x        & = 3
  \end{align*}
  ou seja,
  \[
    \log_2 8 = 3.
  \]
\end{example}
%------------------------------------------------------------------------------%

%------------------------------------------------------------------------------%
\begin{example}
  Encontre \(\log_{\sfrac{1}{81}} 9\)

  \solution
  Mais uma vez, usando a definição de logaritmo
  \begin{align*}
    \log_{\sfrac{1}{81}}9          & = y  \\[2mm]
    \left( \frac{1}{81} \right)^y  & = 9  \\[2mm]
    \left(9^{-2}\right)^y          & = 9  \\[2mm]
    9^{-2y}                        & = 9  \\[2mm]
    -2y                            & = 1  \\[2mm]
    y                              & = -\frac{1}{2}
  \end{align*}
\end{example}
%------------------------------------------------------------------------------%

\emph{Condição de existência do logaritmo:}
A partir da definição de logaritmos, temos que
\[
  a>0 \text{ e } a \neq 1 \text{ e } b > 0.
\]
Portanto, não são definidos, por exemplo, os seguintes logaritmos
\begin{gather*}
  \log_{ 4} (-12) \\
  \log_{ 1}  144  \\
  \log_{ 0}    8  \\
  \log_{-3}   81  \\
  \log_{ 4}    0  \\
  \log_{-6}  -36
\end{gather*}

Note que o valor $\log_b a$ é um expoente.

Segue diretamente da definição que $\log_{a}1=0$,
para todo $a>0$ e $a \neq 1$.

\emph{Logaritmo na base 10}: Quando a base do logaritmo for 10,
podemos omiti-la. $\log 3$ é o logaritmo de $3$ na base $10$.

%------------------------------------------------------------------------------%
\subsection{Gráfico da Função Logarítmica}

Como $a^0 = 1$ e $a^1 = a$, temos que $\log_a 1 = 0$ e $\log_a a = 1$.
Portanto, o gráfico de toda função logarítmica contém os pontos
$(1, 0)$ e $(a, 1)$.
Além disso, sendo a função logarítmica $g(x)=\log_{a}x$ a inversa da
função $g(x)=a^x$, então o gráfico de $f$ é o espelho do gráfico de $g$, com
relação a reta $y=x$.
A Figura~\ref{fig:doislog} ilustra essa situação, mostrando o
esboço do gráfico das funções $g(x)=2^x$ e $\log_{2}x$.

%------------------------------------------------------------------------------%
\begin{figure}
  \centering
  \includegraphics[scale=0.8]{doiselogdedois}
  \caption{Em vermelho, o gráfico da função $g(x)=2^x$; em azul, o gráfico de
    $\log_{2}x$ e em preto tracejado a reta $y=x$..}
  \label{fig:doislog}
\end{figure}
%------------------------------------------------------------------------------%

%------------------------------------------------------------------------------%
\begin{figure}
  \centering
  \includegraphics[width=0.4\textwidth,trim={0 1cm 9cm 1cm},clip]{log1}
  \caption{Gráfico de $f(x)=\log_2x$ (traço contínuo)
    e $g(x)=\log_4x$ (tracejado).}
  \label{fig:grafico_log_crescente}
\end{figure}
%------------------------------------------------------------------------------%

%------------------------------------------------------------------------------%
\begin{figure}
  \centering
  \includegraphics[width=0.4\textwidth,trim={0 1cm 9cm 3cm},clip]{log2}
  \caption{Gráfico de $f(x)=\log_{1/2}x$ (traço contínuo) e
    $g(x)=\log_{1/4}x$ (tracejado).}
  \label{fig:grafico_log_decrescente}
\end{figure}
%------------------------------------------------------------------------------%

A função logarítmica pode ser classificada em crescente ou decrescente.

Se $a>1$ a função $f(x)=\log_a x$ é crescente, pois
\[
  x < y
  \quad\Leftrightarrow\quad
  \log_a x < \log_a y
\]
e o seu gráfico tem o formato das curvas da
Figura~\ref{fig:grafico_log_crescente}.

Se $0<a<1$ a função $f(x)=\log_ax$ é decrescente, pois
\[
  x < y
  \quad\Leftrightarrow\quad
  \log_a x > \log_a y
\]
e o seu gráfico tem o formato das curvas da
Figura~\ref{fig:grafico_log_decrescente}.

Note que o gráfico de $f(x) = \log_{a} x$ não intercepta o eixo $y$, o que
faz sentido, uma vez que o gráfico da sua inversa $y=a^x$ não intercepta o
eixo $x$.

%------------------------------------------------------------------------------%
\emph{Procedimento para esboçar o gráfico de $\log_{a}x$.}

Para esboçar o gráfico de $f(x)=5$, é comum escolhermos 5 valores convenientes
para $x$, a saber as potências de $a$,
\[
  a^{-2}, a^{-1}, a^0, a^1, a^2
\]
e construirmos uma tabela, da forma que fizemos com a função exponencial e, por
conseguinte, marcamos os pontos no gráfico. A escolha desses valores para $x$ é
feita de modo a facilitar o cálculo dos logaritmos.

%------------------------------------------------------------------------------%
\begin{example}
  Esboçar o gráfico da função $f(x)=\log_{3}x$.

  \solution
  Escolhendo
  \[
    x = 3^{-2}, 3^{-1}, 3^0, 3^1, 3^2,
  \]
  podemos montar a tabela
  \begin{ceqn}
  \[
    \begin{array}{l<{\quad}>{\quad}l<{\quad}>{\quad}c}       \hline
      x                     & f(x) = \log_{3} x                 & (x, f(x))           \\ \hline
      \sfrac{1}{9} = 3^{-2} & f(\sfrac{1}{9}) = 3^{-2}          & (\sfrac{1}{9},  -2) \\
      \sfrac{1}{3} = 3^{-1} & f(3) = \log_{3} 3^{-1} = -1       & (\sfrac{1}{3}, -11) \\
      1=3^0                 & f(3) = \log_{3} 1 = 0             & (1, 0)              \\
      3                     & f(3) = \log_{3} 3 = 1             & (3, 1)              \\
      9=3^2                 & f(9) = \log_{3} 9 = \log_3{3^2}=2 & (9, 2)              \\ \hline
    \end{array}
  \]
  \end{ceqn}
  {\large\color{red} Conferir essa tabela!}

  O gráfico de $f$ está na imagem a seguir
  \IncludeGraphicCentredEx{logde3}[0.5]
\end{example}
%------------------------------------------------------------------------------%


